\exercise{Reactivity \label{exMotifReactivity}}{4}
Let $\bf J$ be the Jacobian of a steady state in an ODE system. Consider a small deviation $\vec{\delta}(t)$ from the steady state. Hence the absolute distance of the system from the steady state at time $t$ is $|\vec{\delta}|=\sqrt{\vec{\delta}^{\rm T}\vec{\delta}}$.

\subquestion
Define $a$ as the relative rate of change in $|\vec{\delta}|$, i.e.~the current change in the distance normalized by the current distance, and show 
\eqn{
  a = \frac{1}{|\vec{\delta}|} \frac{{\rm d}}{{\rm d} t} |\vec{\delta}|=\frac{\vec{\delta}^{\dag} {\bf H} \vec{\delta}}{\vec{\delta}^{\dag}\vec{\delta}},
}
where ${\bf H} = ({\bf J}+{\bf J}^{\dag})/2$ is the Hermitian part of the Jacobian. 

\solution 
We know that close to the steady state the deviation changes according to
\eq{
  \dot{\vec{\delta}} = {\bf J} \vec{\delta}.
}
If we transpose and conjugate this equation we find
\eq{
   \dot{\vec{\delta}}^{\dag} = {\vec{\delta}}^{\dag} {\bf J}^{\dag}.
}
The relative change of the absolute distance from the steady state is
\eq{
a = \frac{1}{|\vec{\delta}|} \frac{{\rm d}}{{\rm d} t} |\vec{\delta}|.
}
We now compute
\eqa{
\frac{1}{|\vec{\delta}|} \frac{{\rm d}}{{\rm d} t} |\vec{\delta}| &=& \frac{1}{|\vec{\delta}|} \frac{{\rm d}}{{\rm d} t} \sqrt{\vec{\delta}^\dag\vec{\delta}} \\
&=& \frac{1}{|\vec{\delta}|} \frac{1}{2\sqrt{\vec{\delta}^\dag\vec{\delta}}} \frac{\rm d}{{\rm d} t} \vec{\delta}^\dag\vec{\delta} \\
&=& \frac{1}{2|\vec{\delta}|^2} \left( \dot{\vec{\delta}}^\dag\vec{\delta} + \vec{\delta}^\dag\dot{\vec{\delta}} \right) \\
&=& \frac{1}{2|\vec{\delta}|^2} \left( \vec{\delta}^\dag {\bf J}^\dag \vec{\delta} + \vec{\delta}^\dag{\bf J}\vec{\delta} \right) \\
&=& \frac{\vec{\delta}^\dag {\bf H} \vec{\delta}}{\vec{\delta}^\dag\vec{\delta}},
}
which is the desired result.

\subquestion Define the reactivity of the steady state as the highest possible value of $a$. Show that the reactivity is identical to the rightmost eigenvalue of $\bf H$. 

\solution
Easy, if you see the way: Note that 
\eq{
  a = \frac{\vec{\delta}^{\dag} {\bf H} \vec{\delta}}{\vec{\delta}^{\dag}\vec{\delta}},
}
is the Rayleigh quotient of $\bf H$. Recall that the Rayleigh quotient is maximized if $\vec{\delta}$ is the leading eigenvector and then $a$ is the corresponding eigenvalue. Which shows that the reactivity is the rightmost eigenvalue of $\bf H$.  

\subquestion 
Consider the predator prey system 
\eqan{
\dot{X} &=& X - XY/(3+X) \\
\dot{Y} &=& 2XY/(3+X) - Y.
}
Compute $\bf J$ and $\bf H$ and the reactivity in the nontrivial steady state. 

\solution
The second equation gives us the steady state condition
\eq{
0 = \frac{2XY}{3+X} - Y \\
}
Dividing by $Y$ yields 
\eqa{
0 &=& 2X/(3+X) -1 \\
3+X &=& 2X \\
3 &=& X.
}
So now that we know that $X=3$, we can compute $Y$ from the first row:
\eqa{
 0 &=&  X - XY/(3+X) \\
 0 &=& 3 - 3Y/6\\
 Y &=& 6.
}
So the non-trivial steady state is at $(X^*,Y^*)=(3,6).$ Now we compute 
\eqa{
J_{11} &=& \left.\frac{\partial}{\partial X} \dot{X} \right|_* = \left. 1-\frac{3Y}{(3+X)^2}\right|_* =  1- 18/36 = 0.5  \\
J_{12} &=& \left.\frac{\partial}{\partial Y} \dot{X} \right|_* = \left. - \frac{X}{3+X} \right|_* = -0.5 \\
J_{21} &=& \left.\frac{\partial}{\partial X} \dot{Y} \right|_* = \left.  \frac{6Y}{(3+X)^2}  \right|_*  = 36/36 = 1 \\
J_{22} &=& \left.\frac{\partial}{\partial Y} \dot{Y} \right|_* = \left. \frac{2X}{3+X} - 1 \right|_* = 6/6-1 = 0  
}
Hence the Jacobian is 
\eq{
{\bf J} = \avecc{0.5 & -0.5\\ 1 & 0}  
}
and we can compute 
\eq{
{\bf H} = \frac{\avecc{0.5 & -0.5\\ 1 & 0} + \avecc{0.5 & 1\\ -0.5 & 0}}{2} = \avecc{0.5 & 0.25 \\ 0.25 & 0}
}
Computing the eigenvalues for {\bf H} we get the values $1/4 \pm\sqrt{2}/4$. So the rightmost eigenvalue is $1/4 +\sqrt{2}/4$. Which means the reactivity is $\approx 0.6036$.

\subquestion 
We now consider a larger system where the Jacobian $\bf J$ from (c) appears as a principal submatrix. (Perhaps the system contains $X$ and $Y$ along with many other variables). What can you say about the reactivity of the entire system? 

\solution
For the larger system we compute the reactivity from the Hermitian part of its Jacobian. The Hermitian matrix for the entire system will then contain $\bf H$ from (d) as a principal submatrix. Because the matrix is Hermitian the rightmost eigenvalue of $\bf H$ is a lower bound for the rightmost eigenvalue for the Hermitian matrix for the entire system. Hence the reactivity of the entire system must be at least $1/4+\sqrt{2}/4$.   
